\chapter{Tehnologii folosite}
\section{Spotify API}
API-ul oferit de Spotify a fost utilizat atât pentru construirea setului de date, cât și pentru a putea include 
segmentele similare cu inputul utilizatorului embedded în aplicația Web. Motivul principal pentru care a fost ales
 acest API este functionalitatea de recomandare care permite ca, printr-un request în care este specificat un
 anumit gen muzical, să se obțină o listă de albume și, ulterior, segmente în format audio care se 
încadrează în acel gen. De asemenea, considerând diversitatea genurilor muzicale oferite, 
acest API rezolvă ambele probleme ce au intervenit în obținerea unui set de date.
\section{Librosa}
Librosa este un pachet care pune la dispoziție funcționalități de analiză și prelucrarea a fișierelor audio. 
A fost folosit în această lucrare pentru prelucrarea datelor, extragerea reprezentărilor necesare rețelei neuronale, 
dar și pentru extragerea de informații în vederea curățării setului de date.
\section{Spleeter}
Spleeter\footnote{https://github.com/deezer/spleeter} este o bibliotecă construită cu Tensorflow care folosește modele pre antrenate pentru separarea surselor audio. 
În aceasta lucrare, a fost utilizată pentru a izola elementele vocale de instrumental. O bibliotecă alternativă a fost 
Open-Unmix\footnote{https://github.com/sigsep/open-unmix-pytorch}, dar Spleeter a fost aleasă deoarece are performanțe
 mai bune pe partea de separare vocală\cite{spleeter}.
\section{Tensorflow}
Principalele funcționalități oferite de Tensorflow folosite în această lucrare sunt Keras, necesară pentru construirea 
și antrenarea rețelei neuronale, tf.data API, utilizată pentru crearea unui input pipeline pentru rețeaua construită și 
Tensorboard, necesar pentru vizualizarea unor metrici ale procesului de antrenare.
\section{Sklearn}
Amestecarea datelor, generarea matricelor de confuzie și calcularea celor mai apropiați vectori de predicție se realizează
 utilizând funcționalitățile bibliotecii Sklearn. 
\section{Flask}
Pentru implementarea serverului web responsabil cu primirea fișierelor audio, procesarea lor și generarea de predicții a fost ales framework-ul Flask, în special datorită 
arhitecturii mai simple, în comparație cu alte soluții existente, care se pliază pe nevoile curente ale aplicației.
\section{React}
Pentru implementarea aplicației frontend, a fost utilizat React deoarece facilitează dezvoltarea mai rapidă a aplicației prin utilizarea componentelor 
predefinite și, în comparație cu alternativa Angular, are avantajul simplității.
\section{MongoDB}
Am utilizat un cluster de MongoDB pentru a salva informații adiționale despre melodiile din setul de date precum numele melodiei, albumul 
din care face parte și url-ul de la care poate fi descărcată. Principalul motiv pentru care am optat pentru un serviciu de baze de date în
Cloud este ușurința prin care datele pot fi accesate de pe mai multe dispozitive. Atfel, melodiile au putut fi descărcate în paralel, iar
existența datelor duplicate a fost evitată prin verificarea existenței în baza de date care ține evidența tuturor descărcărilor.